% --- Archon physics formalization source begin ---
% archon:physics
% archon:covers PhyXMiniProblems/problem_phyx_mini_0210.lean
% archon:source-report reports/phyx_mini/problem_phyx_mini_0210.source.json

\chapter{Physics problem phyx\_mini\_0210}
\label{ch:PhyXMiniProblems_problem_phyx_mini_0210}

\paragraph{Problem source.}
\#\# Physical scenario

A highway overpass was observed to resonate as one full loop \textbackslash{}( \textbackslash{}left( \textbackslash{}frac\{1\}\{2\}\textbackslash{}lambda \textbackslash{}right) \textbackslash{}) when a small earthquake shook the ground vertically at \textbackslash{}( 3.0\textbackslash{},\textbackslash{}mathrm\{Hz\} \textbackslash{}). The highway department put a support at the center of the overpass, anchoring it to the ground as shown in figure.

\#\# Question

What resonant frequency would you now expect for the overpass?

\#\# Answer choices

- A: A: 4.0Hz

- B: B: 10.0Hz

- C: C: 8.0Hz

- D: D: 6.0Hz

\#\# Figure caption (auxiliary; use the image as primary evidence)

The image consists of two diagrams labeled "Before modification" and "After modification," showing a side view of a bridge. 

In the "Before modification" diagram:

- The bridge has a flat deck with railings and supports at both ends.
- There is no central support beneath the bridge.
- The background includes clouds and a wavy horizon line.

In the "After modification" diagram:

- A central support has been added under the middle of the bridge marked as "Added support."
- The overall structure of the bridge remains unchanged except for the new support.
- The background is similar, featuring clouds and a wavy horizon.

The main difference between the two diagrams is the added central support in the "After modification" image.

\#\# Dataset metadata

Resonance and Harmonics; Physical Model Grounding Reasoning, Predictive Reasoning

\paragraph{Recorded answer/context.}
D: D: 6.0Hz

\paragraph{Figure/image path.}
/root/proposal\_for\_physic/science-mango/phyx\_mini\_run/phyx\_data/test\_image/210.png

\paragraph{Formalization target.}
create a compiling Lean file with sorry bodies at `PhyXMiniProblems/problem\_phyx\_mini\_0210.lean`. The Lean declarations must preserve the physical quantities, dimensions or dimensional roles, figure labels, governing-law hypotheses, and final relation expressed by this problem.
Use Mathlib/Physlib names found through LeanExplore where available. If a physics API is missing, introduce faithful local abstractions rather than scalar placeholder aliases.

\begin{theorem}[Physics formalization target]
\label{thm:physics:phyx_mini_0210:target}
\lean{PhyXMiniProblems.ProblemPhyXMini0210.problem_phyx_mini_0210}
\uses{def:physics:phyx-mini-0210:phyxminiproblems-problemphyxmini0210-frequencyreadout, def:physics:phyx-mini-0210:phyxminiproblems-problemphyxmini0210-overpassresonancesetup, def:physics:phyx-mini-0210:phyxminiproblems-problemphyxmini0210-matchesproblemdescription, def:physics:phyx-mini-0210:phyxminiproblems-problemphyxmini0210-matchesprimaryfigure, def:physics:phyx-mini-0210:phyxminiproblems-problemphyxmini0210-hasphysicaloverpassparameters, def:physics:phyx-mini-0210:phyxminiproblems-problemphyxmini0210-satisfiesoverpassstandingwavelaws, def:physics:phyx-mini-0210:phyxminiproblems-problemphyxmini0210-iscorrectfrequencyanswer}
The assigned autoformalize agent should translate this physics problem into Lean declarations in the covered file, with theorem and lemma proof bodies written as `by sorry`.
\end{theorem}
\begin{proof}
This is an autoformalization task, not a proof task. Produce statements that can later be proved by the physics prover without weakening the physical model.
\end{proof}

% --- Archon declaration topology begin ---
\section{Declaration topology}
\begin{definition}[Length Quantity]
  \label{def:physics:phyx-mini-0210:phyxminiproblems-problemphyxmini0210-lengthquantity}
  \lean{PhyXMiniProblems.ProblemPhyXMini0210.LengthQuantity}
  A physical longitudinal length or position along the overpass
\end{definition}
\begin{proof}
  This is the declared model definition of Length Quantity.
\end{proof}
\begin{definition}[Frequency Quantity]
  \label{def:physics:phyx-mini-0210:phyxminiproblems-problemphyxmini0210-frequencyquantity}
  \lean{PhyXMiniProblems.ProblemPhyXMini0210.FrequencyQuantity}
  An ordinary physical frequency, with dimension inverse time
\end{definition}
\begin{proof}
  This is the declared model definition of Frequency Quantity.
\end{proof}
\begin{definition}[Speed Quantity]
  \label{def:physics:phyx-mini-0210:phyxminiproblems-problemphyxmini0210-speedquantity}
  \lean{PhyXMiniProblems.ProblemPhyXMini0210.SpeedQuantity}
  A physical wave-propagation speed in the overpass deck
\end{definition}
\begin{proof}
  This is the declared model definition of Speed Quantity.
\end{proof}
\begin{definition}[length Readout]
  \label{def:physics:phyx-mini-0210:phyxminiproblems-problemphyxmini0210-lengthreadout}
  \lean{PhyXMiniProblems.ProblemPhyXMini0210.lengthReadout}
  \uses{def:physics:phyx-mini-0210:phyxminiproblems-problemphyxmini0210-lengthquantity}
  Scalar readout of a physical length in a coherent unit system
\end{definition}
\begin{proof}
  This is the declared model definition of length Readout.
\end{proof}
\begin{definition}[frequency Readout]
  \label{def:physics:phyx-mini-0210:phyxminiproblems-problemphyxmini0210-frequencyreadout}
  \lean{PhyXMiniProblems.ProblemPhyXMini0210.frequencyReadout}
  \uses{def:physics:phyx-mini-0210:phyxminiproblems-problemphyxmini0210-frequencyquantity}
  Scalar readout of an ordinary frequency in a coherent unit system
\end{definition}
\begin{proof}
  This is the declared model definition of frequency Readout.
\end{proof}
\begin{definition}[speed Readout]
  \label{def:physics:phyx-mini-0210:phyxminiproblems-problemphyxmini0210-speedreadout}
  \lean{PhyXMiniProblems.ProblemPhyXMini0210.speedReadout}
  \uses{def:physics:phyx-mini-0210:phyxminiproblems-problemphyxmini0210-speedquantity}
  Scalar readout of Physlib's nonnegative speed quantity
\end{definition}
\begin{proof}
  This is the declared model definition of speed Readout.
\end{proof}
\begin{definition}[Overpass Configuration]
  \label{def:physics:phyx-mini-0210:phyxminiproblems-problemphyxmini0210-overpassconfiguration}
  \lean{PhyXMiniProblems.ProblemPhyXMini0210.OverpassConfiguration}
  The two panels explicitly labeled in the supplied figure
\end{definition}
\begin{proof}
  This is the declared model definition of Overpass Configuration.
\end{proof}
\begin{definition}[Support Site]
  \label{def:physics:phyx-mini-0210:phyxminiproblems-problemphyxmini0210-supportsite}
  \lean{PhyXMiniProblems.ProblemPhyXMini0210.SupportSite}
  Longitudinal sites at which the figure may show a support
\end{definition}
\begin{proof}
  This is the declared model definition of Support Site.
\end{proof}
\begin{definition}[Support Kind]
  \label{def:physics:phyx-mini-0210:phyxminiproblems-problemphyxmini0210-supportkind}
  \lean{PhyXMiniProblems.ProblemPhyXMini0210.SupportKind}
  Mechanical support visible at a site in a given panel
\end{definition}
\begin{proof}
  This is the declared model definition of Support Kind.
\end{proof}
\begin{definition}[Transverse Boundary Condition]
  \label{def:physics:phyx-mini-0210:phyxminiproblems-problemphyxmini0210-transverseboundarycondition}
  \lean{PhyXMiniProblems.ProblemPhyXMini0210.TransverseBoundaryCondition}
  Transverse-displacement condition of a standing mode at a marked site
\end{definition}
\begin{proof}
  This is the declared model definition of Transverse Boundary Condition.
\end{proof}
\begin{definition}[Ground Motion Direction]
  \label{def:physics:phyx-mini-0210:phyxminiproblems-problemphyxmini0210-groundmotiondirection}
  \lean{PhyXMiniProblems.ProblemPhyXMini0210.GroundMotionDirection}
  Direction of the earthquake's ground displacement
\end{definition}
\begin{proof}
  This is the declared model definition of Ground Motion Direction.
\end{proof}
\begin{definition}[Excitation Source]
  \label{def:physics:phyx-mini-0210:phyxminiproblems-problemphyxmini0210-excitationsource}
  \lean{PhyXMiniProblems.ProblemPhyXMini0210.ExcitationSource}
  Qualitative source of the external drive in the problem statement
\end{definition}
\begin{proof}
  This is the declared model definition of Excitation Source.
\end{proof}
\begin{definition}[Figure Label]
  \label{def:physics:phyx-mini-0210:phyxminiproblems-problemphyxmini0210-figurelabel}
  \lean{PhyXMiniProblems.ProblemPhyXMini0210.FigureLabel}
  Text labels printed on the two panels and the new support
\end{definition}
\begin{proof}
  This is the declared model definition of Figure Label.
\end{proof}
\begin{definition}[Figure Element]
  \label{def:physics:phyx-mini-0210:phyxminiproblems-problemphyxmini0210-figureelement}
  \lean{PhyXMiniProblems.ProblemPhyXMini0210.FigureElement}
  Figure elements to which the printed labels point
\end{definition}
\begin{proof}
  This is the declared model definition of Figure Element.
\end{proof}
\begin{definition}[Overpass Resonance Setup]
  \label{def:physics:phyx-mini-0210:phyxminiproblems-problemphyxmini0210-overpassresonancesetup}
  \lean{PhyXMiniProblems.ProblemPhyXMini0210.OverpassResonanceSetup}
  \uses{def:physics:phyx-mini-0210:phyxminiproblems-problemphyxmini0210-lengthquantity, def:physics:phyx-mini-0210:phyxminiproblems-problemphyxmini0210-frequencyquantity, def:physics:phyx-mini-0210:phyxminiproblems-problemphyxmini0210-speedquantity, def:physics:phyx-mini-0210:phyxminiproblems-problemphyxmini0210-overpassconfiguration, def:physics:phyx-mini-0210:phyxminiproblems-problemphyxmini0210-supportsite, def:physics:phyx-mini-0210:phyxminiproblems-problemphyxmini0210-supportkind, def:physics:phyx-mini-0210:phyxminiproblems-problemphyxmini0210-transverseboundarycondition, def:physics:phyx-mini-0210:phyxminiproblems-problemphyxmini0210-groundmotiondirection, def:physics:phyx-mini-0210:phyxminiproblems-problemphyxmini0210-excitationsource, def:physics:phyx-mini-0210:phyxminiproblems-problemphyxmini0210-figurelabel, def:physics:phyx-mini-0210:phyxminiproblems-problemphyxmini0210-figureelement}
  Physical quantities and figure-indexed data for the overpass. halfWavelengthLoopCount counts node-to-node loops across the entire deck. The resonant frequency fields are unconstrained here: in particular, the post-modification frequency is not built into the setup
\end{definition}
\begin{proof}
  This is the declared model definition of Overpass Resonance Setup.
\end{proof}
\begin{definition}[Matches Problem Description]
  \label{def:physics:phyx-mini-0210:phyxminiproblems-problemphyxmini0210-matchesproblemdescription}
  \lean{PhyXMiniProblems.ProblemPhyXMini0210.MatchesProblemDescription}
  \uses{def:physics:phyx-mini-0210:phyxminiproblems-problemphyxmini0210-frequencyreadout, def:physics:phyx-mini-0210:phyxminiproblems-problemphyxmini0210-overpassresonancesetup}
  Data stated in the prose: the drive is a small earthquake with vertical ground motion, the observed initial mode is one half-wavelength loop at 3.0 Hz, and the requested resonance is the one after modification
\end{definition}
\begin{proof}
  This is the declared model definition of Matches Problem Description.
\end{proof}
\begin{definition}[Matches Primary Figure]
  \label{def:physics:phyx-mini-0210:phyxminiproblems-problemphyxmini0210-matchesprimaryfigure}
  \lean{PhyXMiniProblems.ProblemPhyXMini0210.MatchesPrimaryFigure}
  \uses{def:physics:phyx-mini-0210:phyxminiproblems-problemphyxmini0210-lengthreadout, def:physics:phyx-mini-0210:phyxminiproblems-problemphyxmini0210-overpassresonancesetup}
  Readouts from the supplied image. Both panels show the same deck and its two end supports. The center site is empty before modification and is a grounded support afterward. Its longitudinal coordinate is exactly the midpoint, so the support divides the deck into two intervals. The three printed labels are also recorded explicitly
\end{definition}
\begin{proof}
  This is the declared model definition of Matches Primary Figure.
\end{proof}
\begin{definition}[Has Physical Overpass Parameters]
  \label{def:physics:phyx-mini-0210:phyxminiproblems-problemphyxmini0210-hasphysicaloverpassparameters}
  \lean{PhyXMiniProblems.ProblemPhyXMini0210.HasPhysicalOverpassParameters}
  \uses{def:physics:phyx-mini-0210:phyxminiproblems-problemphyxmini0210-lengthreadout, def:physics:phyx-mini-0210:phyxminiproblems-problemphyxmini0210-frequencyreadout, def:physics:phyx-mini-0210:phyxminiproblems-problemphyxmini0210-speedreadout, def:physics:phyx-mini-0210:phyxminiproblems-problemphyxmini0210-overpassresonancesetup}
  Positivity conditions selecting nondegenerate overpass modes
\end{definition}
\begin{proof}
  This is the declared model definition of Has Physical Overpass Parameters.
\end{proof}
\begin{definition}[Satisfies Overpass Standing Wave Laws]
  \label{def:physics:phyx-mini-0210:phyxminiproblems-problemphyxmini0210-satisfiesoverpassstandingwavelaws}
  \lean{PhyXMiniProblems.ProblemPhyXMini0210.SatisfiesOverpassStandingWaveLaws}
  \uses{def:physics:phyx-mini-0210:phyxminiproblems-problemphyxmini0210-lengthreadout, def:physics:phyx-mini-0210:phyxminiproblems-problemphyxmini0210-frequencyreadout, def:physics:phyx-mini-0210:phyxminiproblems-problemphyxmini0210-speedreadout, def:physics:phyx-mini-0210:phyxminiproblems-problemphyxmini0210-overpassresonancesetup}
  Standing-wave laws for the lowest mode compatible with the supports: every mechanical support anchors a transverse-displacement node; the lowest mode has one half-wavelength loop on each support interval; n λ = 2 L for n loops across a deck of length L; wave kinematics gives v = f λ; adding the ideal support does not change the deck's propagation speed. These are generic physical/modeling relations. None states the requested post-modification frequency or names an answer choice
\end{definition}
\begin{proof}
  This is the declared model definition of Satisfies Overpass Standing Wave Laws.
\end{proof}
\begin{lemma}[post Modification Loop Count eq two]
  \label{lem:physics:phyx-mini-0210:phyxminiproblems-problemphyxmini0210-postmodificationloopcount-eq-two}
  \lean{PhyXMiniProblems.ProblemPhyXMini0210.postModificationLoopCount_eq_two}
  \uses{def:physics:phyx-mini-0210:phyxminiproblems-problemphyxmini0210-overpassresonancesetup, def:physics:phyx-mini-0210:phyxminiproblems-problemphyxmini0210-matchesprimaryfigure, def:physics:phyx-mini-0210:phyxminiproblems-problemphyxmini0210-satisfiesoverpassstandingwavelaws}
  The added midpoint node makes the lowest compatible mode contain two loops, one on each half of the overpass
\end{lemma}
\begin{proof}
  The conclusion follows from the governing relations and figure readouts named in the statement of post Modification Loop Count eq two.
\end{proof}
\begin{lemma}[twice post Modification Wavelength eq before Modification Wavelength]
  \label{lem:physics:phyx-mini-0210:phyxminiproblems-problemphyxmini0210-twice-postmodificationwavelength-eq-beforemodificationwavelength}
  \lean{PhyXMiniProblems.ProblemPhyXMini0210.twice_postModificationWavelength_eq_beforeModificationWavelength}
  \uses{def:physics:phyx-mini-0210:phyxminiproblems-problemphyxmini0210-lengthreadout, def:physics:phyx-mini-0210:phyxminiproblems-problemphyxmini0210-overpassresonancesetup, def:physics:phyx-mini-0210:phyxminiproblems-problemphyxmini0210-matchesproblemdescription, def:physics:phyx-mini-0210:phyxminiproblems-problemphyxmini0210-matchesprimaryfigure, def:physics:phyx-mini-0210:phyxminiproblems-problemphyxmini0210-hasphysicaloverpassparameters, def:physics:phyx-mini-0210:phyxminiproblems-problemphyxmini0210-satisfiesoverpassstandingwavelaws}
  The wavelength of the new lowest mode is half the wavelength of the observed one-loop mode before modification
\end{lemma}
\begin{proof}
  The conclusion follows from the governing relations and figure readouts named in the statement of twice post Modification Wavelength eq before Modification Wavelength.
\end{proof}
\begin{lemma}[post Modification Frequency eq twice before Modification Frequency]
  \label{lem:physics:phyx-mini-0210:phyxminiproblems-problemphyxmini0210-postmodificationfrequency-eq-twice-beforemodificationfrequency}
  \lean{PhyXMiniProblems.ProblemPhyXMini0210.postModificationFrequency_eq_twice_beforeModificationFrequency}
  \uses{def:physics:phyx-mini-0210:phyxminiproblems-problemphyxmini0210-frequencyreadout, def:physics:phyx-mini-0210:phyxminiproblems-problemphyxmini0210-overpassresonancesetup, def:physics:phyx-mini-0210:phyxminiproblems-problemphyxmini0210-matchesproblemdescription, def:physics:phyx-mini-0210:phyxminiproblems-problemphyxmini0210-matchesprimaryfigure, def:physics:phyx-mini-0210:phyxminiproblems-problemphyxmini0210-hasphysicaloverpassparameters, def:physics:phyx-mini-0210:phyxminiproblems-problemphyxmini0210-satisfiesoverpassstandingwavelaws}
  With unchanged wave speed and v = f λ, halving the wavelength doubles the ordinary resonant frequency
\end{lemma}
\begin{proof}
  The conclusion follows from the governing relations and figure readouts named in the statement of post Modification Frequency eq twice before Modification Frequency.
\end{proof}
\begin{definition}[Answer Choice]
  \label{def:physics:phyx-mini-0210:phyxminiproblems-problemphyxmini0210-answerchoice}
  \lean{PhyXMiniProblems.ProblemPhyXMini0210.AnswerChoice}
  Labels of the four answer choices printed in the problem
\end{definition}
\begin{proof}
  This is the declared model definition of Answer Choice.
\end{proof}
\begin{definition}[frequency Hz]
  \label{def:physics:phyx-mini-0210:phyxminiproblems-problemphyxmini0210-answerchoice-frequencyhz}
  \lean{PhyXMiniProblems.ProblemPhyXMini0210.AnswerChoice.frequencyHz}
  \uses{def:physics:phyx-mini-0210:phyxminiproblems-problemphyxmini0210-answerchoice}
  Frequency in hertz printed beside each answer choice
\end{definition}
\begin{proof}
  This is the declared model definition of frequency Hz.
\end{proof}
\begin{definition}[Is Correct Frequency Answer]
  \label{def:physics:phyx-mini-0210:phyxminiproblems-problemphyxmini0210-iscorrectfrequencyanswer}
  \lean{PhyXMiniProblems.ProblemPhyXMini0210.IsCorrectFrequencyAnswer}
  \uses{def:physics:phyx-mini-0210:phyxminiproblems-problemphyxmini0210-frequencyreadout, def:physics:phyx-mini-0210:phyxminiproblems-problemphyxmini0210-overpassresonancesetup, def:physics:phyx-mini-0210:phyxminiproblems-problemphyxmini0210-answerchoice}
  An answer choice matches the requested post-modification SI frequency
\end{definition}
\begin{proof}
  This is the declared model definition of Is Correct Frequency Answer.
\end{proof}
% --- Archon declaration topology end ---
% --- Archon physics formalization source end ---
